\section{Experimental Results and Analysis}
This section presents the comprehensive experimental results obtained through our nested cross-validation framework with Optuna hyperparameter optimization. All experiments were conducted using the rigorous methodology described in previous sections, ensuring reliable and statistically significant performance estimates.

\subsection{Experimental Setup}

\subsubsection{Training Configuration}

All experiments were conducted with the following standardized configuration to ensure fair comparison across different dataset variants and model architectures:

\begin{itemize}
    \item \textbf{Training Epochs}: 50 epochs for YOLOv5s and YOLOv8s comparison studies, 5 epochs for additional experimental variants
    \item \textbf{Cross-Validation}: 5 inner folds for hyperparameter optimization, 5 outer folds for final evaluation
    \item \textbf{Optuna Trials}: 20 hyperparameter optimization trials per inner fold
    \item \textbf{Hardware}: Dual NVIDIA A100 GPU training configuration (GPU IDs 6 and 7) with 40GB memory per GPU
    \item \textbf{Evaluation Metric}: mAP50-95 (mean Average Precision at IoU thresholds 0.5 to 0.95)
\end{itemize}

\subsubsection{Command Execution}

The standardized training command used across all experiments was:

\begin{lstlisting}[language=bash]
# Standard training command template
python main.py \
  --model [MODEL] \
  --data_path ./pipeline_dataset_[VARIANT] \
  --gpu_ids "6,7" \
  --nested_cv \
  --n_trials 20 \
  --epochs [EPOCHS] \
  --n_folds 5 \
  --n_outer_folds 5
\end{lstlisting}

where \texttt{[MODEL]} represents the YOLO variant and \texttt{[VARIANT]} represents the specific dataset variant.

\subsubsection{Hardware Specifications}

All experiments were conducted on a high-performance computing system equipped with dual NVIDIA A100 GPUs, providing the computational power necessary for extensive agricultural computer vision research:

\begin{itemize}
    \item \textbf{GPU Architecture}: Two NVIDIA A100 Tensor Core GPUs (40GB HBM2e memory per GPU)
    \item \textbf{Total GPU Memory}: 80GB combined GPU memory for handling high-resolution agricultural imagery
    \item \textbf{CUDA Capability}: 8.0, enabling advanced tensor operations and mixed-precision training
    \item \textbf{Memory Bandwidth}: 1,555 GB/s per GPU, crucial for efficient data transfer during intensive training
    \item \textbf{Tensor Performance}: 312 TFLOPS (Tensor Float 32) for accelerated deep learning computations
    \item \textbf{Configuration}: GPUs utilized with IDs 6 and 7 in distributed training mode
\end{itemize}

The A100 architecture is particularly well-suited for agricultural computer vision tasks due to its high memory capacity (essential for processing high-resolution field images), advanced tensor operations (optimizing YOLO model training), and multi-instance GPU capabilities (enabling efficient hyperparameter optimization across multiple trials simultaneously).

\subsection{Dataset Variants Evaluated}

Our comprehensive evaluation covers the four primary dataset variants, each representing different levels of classification complexity:

\begin{table}[H]
\centering
\caption{Evaluated Dataset Variants}
\label{tab:evaluated_variants}
\begin{tabular}{@{}lccp{6cm}@{}}
\toprule
\textbf{Dataset Variant} & \textbf{Classes} & \textbf{Complexity} & \textbf{Application Focus} \\
\midrule
Fine24 & 24 & High & Detailed species-level classification \\
CropsOrWeed9 & 9 & Medium & Crop type identification with weed detection \\
CropOrWeed2 & 2 & Low & Basic crop vs. weed distinction \\
Coarse1 & 1 & Minimal & Vegetation detection and localization \\
\bottomrule
\end{tabular}
\end{table}

\subsubsection{Understanding the Experimental Dataset Variants}

This table presents the strategic selection of dataset variants used in our comprehensive evaluation, representing a carefully designed spectrum of agricultural computer vision challenges. Understanding these variants is crucial for interpreting our results and applying them to real-world agricultural scenarios.

\textbf{What This Table Contains:} The table organizes our experimental variants by complexity level, from the most challenging (Fine24) to the simplest (Coarse1). Each variant represents a different agricultural use case, with the "Classes" column indicating the number of distinct plant categories the system must distinguish, the "Complexity" column providing a qualitative assessment of the detection challenge, and the "Application Focus" column describing the primary agricultural purpose.

\textbf{Strategic Variant Selection and Their Practical Significance:}

\begin{itemize}
    \item \textbf{Fine24 (24 classes, High complexity)}: This represents the most ambitious agricultural computer vision challenge, requiring the system to distinguish between 8 different crop types and 16 specific weed species. This level of detail is invaluable for advanced precision agriculture, biodiversity research, and comprehensive farm management systems. The high complexity serves as a stress test for our methodology—if the approach works well here, it demonstrates robustness for any agricultural application. For practical deployment, this variant would support applications like detailed crop health monitoring, species-specific treatment recommendations, and regulatory compliance for organic farming where precise weed identification is crucial.
    
    \item \textbf{CropsOrWeed9 (9 classes, Medium complexity)}: This variant strikes a balance between practical utility and manageable complexity, distinguishing between 8 crop types plus a general weed category. This configuration is particularly valuable for mixed farming operations, crop rotation planning, and agricultural management systems that need to track different crop types across seasons. The medium complexity makes it suitable for commercial deployment while providing sufficient detail for meaningful farm management decisions. In practical terms, this variant could support applications like automated crop mapping, yield prediction systems, and multi-crop precision agriculture tools.
    
    \item \textbf{CropOrWeed2 (2 classes, Low complexity)}: This binary classification addresses the fundamental agricultural question: "Is this a plant I want to keep (crop) or remove (weed)?" This simplicity makes it ideal for the most immediate practical applications in precision agriculture, such as automated weeding systems, targeted herbicide application, and basic crop monitoring. The low complexity increases the likelihood of reliable performance in real-world deployment, making it the most accessible entry point for farmers adopting computer vision technology. This variant directly supports the core economic goal of agriculture: maximizing useful plant growth while minimizing harmful plant competition.
    
    \item \textbf{Coarse1 (1 class, Minimal complexity)}: This single-class vegetation detection focuses on the basic task of identifying "where plants are growing" regardless of their type. While seemingly simple, this variant is crucial for applications like field coverage assessment, irrigation planning, bare soil identification, and general crop monitoring. The minimal complexity provides a foundation for more complex systems and serves as a baseline for understanding fundamental computer vision performance in agricultural environments. For practical deployment, this variant supports large-scale field surveying, growth monitoring over time, and integration with precision agriculture equipment that needs basic vegetation maps.
\end{itemize}

\textbf{Experimental Design Strategy:} The selection of these four variants follows a deliberate experimental design that enables several critical insights:

\begin{itemize}
    \item \textbf{Scalability Assessment}: By testing across this complexity spectrum, we can evaluate how well our methodology scales from simple to challenging agricultural tasks, providing guidance for selecting appropriate complexity levels for different applications.
    
    \item \textbf{Performance Trade-offs}: The variants reveal the relationship between task complexity and achievable performance, helping agricultural technology developers make informed decisions about feature complexity versus reliability.
    
    \item \textbf{Deployment Readiness}: Each variant represents a different level of deployment readiness, from research-grade (Fine24) to commercially viable (CropOrWeed2) to foundational infrastructure (Coarse1).
    
    \item \textbf{Economic Impact Assessment}: The variants correspond to different economic value propositions, allowing evaluation of which complexity levels provide the best return on investment for agricultural technology deployment.
\end{itemize}

\textbf{Practical Application Guidance:} This variant structure provides clear guidance for agricultural stakeholders: research institutions and advanced operations might focus on Fine24 for comprehensive analysis, commercial farms with diverse crops would benefit from CropsOrWeed9, precision agriculture implementers should prioritize CropOrWeed2 for immediate impact, and large-scale operations might start with Coarse1 for basic field monitoring infrastructure.

\subsection{Comprehensive Experimental Results}

\subsubsection{YOLOv8s Performance Results - 50 Epochs}

The following table presents the comprehensive results for YOLOv8s trained for 50 epochs across all dataset variants:

\begin{table}[H]
\centering
\caption{YOLOv8s Performance Results - 50 Epochs (mAP50-95)}
\label{tab:yolov8s_50_results}
\resizebox{\textwidth}{!}{%
\begin{tabular}{@{}lccccccc@{}}
\toprule
\textbf{Dataset} & \textbf{Fold 1} & \textbf{Fold 2} & \textbf{Fold 3} & \textbf{Fold 4} & \textbf{Fold 5} & \textbf{Mean} & \textbf{Std Dev} \\
\midrule
Fine24 & 0.461 & 0.468 & 0.460 & 0.457 & 0.449 & \textbf{0.459} & ±0.006 \\
CropsOrWeed9 & 0.647 & 0.705 & 0.691 & 0.682 & 0.685 & \textbf{0.682} & ±0.019 \\
CropOrWeed2 & 0.578 & 0.569 & 0.559 & 0.540 & 0.571 & \textbf{0.563} & ±0.013 \\
Coarse1 & 0.428 & 0.442 & 0.433 & 0.420 & 0.427 & \textbf{0.430} & ±0.007 \\
\bottomrule
\end{tabular}%
}
\end{table}

\subsubsection{Comprehensive Interpretation of YOLOv8s Performance Results}

This table presents the most important findings of our research: the actual performance achieved by the YOLOv8s model across different agricultural detection tasks. Understanding these results is crucial for making informed decisions about deploying computer vision technology in agricultural settings.

\textbf{What This Table Contains:} The table shows performance results using the mAP50-95 metric (mean Average Precision at Intersection over Union thresholds from 0.5 to 0.95), which is the gold standard for evaluating object detection systems. Each row represents a different agricultural task complexity, and each column (Fold 1-5) represents an independent test on different portions of the data to ensure reliability. The numbers range from 0 to 1, where 1 represents perfect detection and 0 represents complete failure.

\textbf{Detailed Performance Analysis by Agricultural Application:}

\begin{itemize}
    \item \textbf{CropsOrWeed9 (0.682 mean performance)}: This represents our best-performing configuration and is particularly significant for practical agricultural applications. With 68.2\% accuracy in distinguishing between 8 different crop types plus weeds, this system could reliably support farm management decisions. The relatively low standard deviation (±0.019) indicates consistent performance across different field conditions. For practical deployment, this means farmers could trust the system to correctly identify crop types about 7 times out of 10, which is sufficient for applications like crop mapping, yield estimation, and general field monitoring.
    
    \item \textbf{CropOrWeed2 (0.563 mean performance)}: With 56.3\% accuracy for the fundamental agricultural question "crop or weed?", this system shows promise for precision agriculture applications. While not perfect, this level of performance could support semi-automated systems where farmers review recommendations before action. The standard deviation of ±0.013 shows good consistency. In practical terms, this system could guide targeted treatments while requiring human oversight for critical decisions that could damage crops.
    
    \item \textbf{Fine24 (0.459 mean performance)}: At 45.9\% accuracy for detailed species-level identification across 24 different plant types, this represents the current state-of-the-art for complex agricultural computer vision. While the accuracy is moderate, the very low standard deviation (±0.006) indicates that this performance level is highly reliable and predictable. This system would be valuable for research applications, biodiversity studies, and specialized farming operations where detailed species information is crucial, particularly when combined with human expertise.
    
    \item \textbf{Coarse1 (0.430 mean performance)}: The 43.0\% performance for basic vegetation detection seems surprisingly low for such a simple task, but this reflects the challenging nature of agricultural environments where vegetation boundaries are often unclear due to varying lighting, overlapping plants, and complex backgrounds. The low standard deviation (±0.007) indicates consistent performance. This system could serve as a foundation for more complex applications or as a first-stage filter in multi-step detection systems.
\end{itemize}

\textbf{Reliability and Deployment Readiness:} The consistently low standard deviations across all variants (ranging from ±0.006 to ±0.019) demonstrate that our methodology produces reliable, predictable results. This reliability is crucial for agricultural deployment, where inconsistent performance could lead to crop damage or economic losses. The five-fold validation ensures that these performance estimates are trustworthy and not dependent on specific field conditions or image characteristics.

\textbf{Practical Implications for Farmers and Agricultural Technology Developers:} These results suggest that computer vision technology is ready for supportive roles in agriculture, particularly for crop type identification (CropsOrWeed9) and general crop-weed distinction (CropOrWeed2). However, the technology should be deployed as decision support tools rather than fully autonomous systems, allowing farmers to review and approve recommendations before implementation.

\subsubsection{YOLOv5su Performance Results (50 Epochs)}

Comparative results for YOLOv5su trained for 50 epochs:

\begin{table}[H]
\centering
\caption{YOLOv5su Performance Results - 50 Epochs (mAP50-95)}
\label{tab:yolov5su_50_results}
\resizebox{\textwidth}{!}{%
\begin{tabular}{@{}lccccccc@{}}
\toprule
\textbf{Dataset} & \textbf{Fold 1} & \textbf{Fold 2} & \textbf{Fold 3} & \textbf{Fold 4} & \textbf{Fold 5} & \textbf{Mean} & \textbf{Std Dev} \\
\midrule
Fine24 & 0.462 & 0.469 & 0.455 & 0.458 & 0.449 & \textbf{0.459} & ±0.007 \\
CropsOrWeed9 & 0.674 & 0.686 & 0.690 & 0.671 & 0.685 & \textbf{0.681} & ±0.008 \\
CropOrWeed2 & 0.570 & 0.566 & 0.545 & 0.545 & 0.563 & \textbf{0.558} & ±0.011 \\
Coarse1 & 0.427 & 0.437 & 0.430 & 0.420 & 0.423 & \textbf{0.427} & ±0.006 \\
\bottomrule
\end{tabular}%
}
\end{table}

\subsubsection{Comprehensive Interpretation of YOLOv5su Performance Results}

This table presents the performance results for the YOLOv5su architecture, providing a critical comparison point to understand how different YOLO variants perform on agricultural computer vision tasks. The YOLOv5su model represents a mature, well-established architecture that has been widely adopted in computer vision applications, making it an important benchmark for agricultural deployment considerations.

\textbf{What This Table Contains:} Like the YOLOv8s results, this table shows mAP50-95 performance metrics across five cross-validation folds for each dataset variant. The consistent formatting allows for direct comparison between architectures, enabling informed decisions about which model variant offers the best trade-offs for specific agricultural applications.

\textbf{Detailed Performance Analysis by Agricultural Application:}

\begin{itemize}
    \item \textbf{CropsOrWeed9 (0.681 mean performance)}: YOLOv5su achieves 68.1\% accuracy for crop type identification, which is remarkably close to the YOLOv8s performance (68.2\%). This near-identical performance (difference of only 0.1\%) is significant because it demonstrates that both modern and established YOLO architectures can achieve practical-level performance for agricultural applications. The low standard deviation (±0.008) indicates excellent reliability. For farmers and agricultural technology developers, this means they can choose between architectures based on factors like deployment compatibility, computational requirements, or existing infrastructure rather than worrying about significant performance differences.
    
    \item \textbf{CropOrWeed2 (0.558 mean performance)}: With 55.8\% accuracy for basic crop-versus-weed detection, YOLOv5su performs slightly lower than YOLOv8s (56.3\%) but the difference is practically negligible (0.5\%). This performance level still supports decision-support applications for precision agriculture. The standard deviation of ±0.011 shows good consistency, though slightly higher variability than YOLOv8s. In practical terms, both architectures offer similar reliability for farmers implementing semi-automated weed management systems.
    
    \item \textbf{Fine24 (0.459 mean performance)}: Remarkably, YOLOv5su achieves identical performance to YOLOv8s (45.9\%) for the most complex 24-class species identification task. This identical performance across architectures is particularly noteworthy because it suggests that the limiting factor for complex agricultural computer vision is not the model architecture but rather the inherent difficulty of the task and the quality of available training data. The slightly higher standard deviation (±0.007 vs ±0.006) indicates marginally less consistency but still excellent reliability for research and specialized agricultural applications.
    
    \item \textbf{Coarse1 (0.427 mean performance)}: For basic vegetation detection, YOLOv5su achieves 42.7\% accuracy, slightly lower than YOLOv8s (43.0\%). While this represents the lowest performance across all variants, the very low standard deviation (±0.006) demonstrates consistent, predictable behavior. This consistency is valuable for applications where reliable baseline performance is more important than peak accuracy, such as large-scale field monitoring or as a first-stage filter in multi-step detection systems.
\end{itemize}

\textbf{Architecture Comparison Insights:} The YOLOv5su results reveal several important insights about agricultural computer vision deployment:

\begin{itemize}
    \item \textbf{Performance Parity}: The minimal differences between YOLOv5su and YOLOv8s (ranging from 0.0\% to 0.9\%) suggest that architecture choice should be based on practical deployment considerations rather than pure performance metrics.
    
    \item \textbf{Consistency Validation}: Both architectures show similar patterns of performance across complexity levels, validating that our methodology successfully optimizes different models to their potential rather than favoring one particular architecture.
    
    \item \textbf{Deployment Flexibility}: The comparable performance means agricultural technology developers can choose the architecture that best fits their existing systems, computational constraints, or support infrastructure without sacrificing detection accuracy.
\end{itemize}

\textbf{Practical Implications for Agricultural Deployment:} These results demonstrate that YOLOv5su represents a viable, mature option for agricultural computer vision deployment. Organizations with existing YOLOv5 infrastructure, computational constraints that favor the older architecture, or requirements for well-established, thoroughly tested systems can confidently deploy YOLOv5su-based solutions knowing they will achieve performance levels comparable to the newest architectures. This flexibility reduces the technological barriers to agricultural computer vision adoption and enables more diverse deployment scenarios.

\subsubsection{Performance Comparison - 50 Epochs}

Our results demonstrate significant improvements compared to the original CropAndWeed paper benchmarks:

\begin{table}[H]
\centering
\caption{Comparison with Original CropAndWeed Paper Results (50 Epochs)}
\label{tab:baseline_comparison_50}
\resizebox{\textwidth}{!}{%
\begin{tabular}{@{}lcccl@{}}
\toprule
\textbf{Model/Dataset} & \textbf{Our Results} & \textbf{Original} & \textbf{Improvement} & \textbf{Method} \\
\midrule
YOLOv8s Fine24 & \textbf{0.459} & 0.331 & +38.67\% & Nested CV \\
YOLOv8s CropsOrWeed9 & \textbf{0.682} & 0.322 & +111.80\% & + Optuna \\
YOLOv8s CropOrWeed2 & \textbf{0.563} & 0.330 & +70.61\% & \\
YOLOv8s Coarse1 & \textbf{0.430} & 0.297 & +44.78\%  \\
\midrule
YOLOv5su Fine24 & \textbf{0.459} & 0.331 & +38.67\% & Nested CV \\
YOLOv5su CropsOrWeed9 & \textbf{0.681} & 0.322 & +111.49\% & + Optuna \\
YOLOv5su CropOrWeed2 & \textbf{0.558} & 0.330 & +69.09\% &  \\
YOLOv5su Coarse1 & \textbf{0.427} & 0.297 & +43.77\%  \\
\bottomrule
\end{tabular}%
}
\end{table}

\subsubsection{Understanding the Magnitude of Methodological Improvements}

This comparison table represents one of the most significant findings of our research: the dramatic impact that proper methodology can have on agricultural computer vision performance. These improvements demonstrate that the technology is advancing much more rapidly than might be apparent from traditional research approaches.

\textbf{What This Table Reveals:} This table compares our results (using nested cross-validation and Optuna hyperparameter optimization) against the original CropAndWeed paper results using the same models and datasets. The "Improvement" column shows percentage increases in performance, representing real advances in agricultural detection capability that translate directly to better farming outcomes.

\textbf{Breakthrough Results Analysis:}

\begin{itemize}
    \item \textbf{CropsOrWeed9 - Over 111\% Improvement}: The most remarkable finding is the more-than-doubling of performance for crop type identification (from 32.2\% to 68.2\% accuracy). This improvement transforms the technology from a research curiosity to a practically deployable tool. In concrete terms, this means a system that previously failed more than two-thirds of the time now succeeds more than two-thirds of the time. For farmers, this represents the difference between technology that cannot be trusted for important decisions and technology that can serve as a reliable decision support tool.
    
    \item \textbf{CropOrWeed2 - 70\% Improvement}: The 70\% improvement in basic crop-versus-weed detection (from 33.0\% to 56.3\%) represents crossing a critical threshold from unreliable to potentially useful. While still requiring human oversight, this performance level could support semi-automated systems for targeted herbicide application or mechanical weeding guidance. The improvement represents progress from a system that farmers would reject as useless to one they might consider as an assistant.
    
    \item \textbf{Fine24 - 38\% Improvement}: Even for the most complex task (24-class species identification), the 38\% improvement (from 33.1\% to 45.9\%) demonstrates that sophisticated agricultural computer vision is advancing rapidly. This level of performance makes the technology valuable for research applications and specialized farming operations where detailed species information justifies the current limitations.
    
    \item \textbf{Coarse1 - 44\% Improvement}: The 44\% improvement in basic vegetation detection (from 29.7\% to 43.0\%) shows that even apparently simple agricultural computer vision tasks benefit significantly from proper methodology. This improvement enables more reliable field monitoring and coverage assessment applications.
\end{itemize}

\textbf{Methodological Significance:} The consistency of improvements across all dataset variants and model architectures (ranging from 38\% to 111\%) demonstrates that these gains are not due to luck or overfitting, but represent genuine methodological advances. The "Method" column highlights that these improvements come from nested cross-validation and Optuna optimization—techniques that can be applied to any agricultural computer vision project.

\textbf{Practical Impact for Agricultural Technology Development:} These results suggest that many agricultural computer vision systems deployed today could be dramatically improved simply by applying better training and evaluation methodologies. For agricultural technology companies, this represents an opportunity to significantly enhance existing products without developing new algorithms. For farmers, this means that computer vision technology may be more mature and ready for deployment than previously thought.

\textbf{Economic and Environmental Implications:} The improvements shown in this table translate directly to economic and environmental benefits. More accurate crop-weed detection enables reduced herbicide usage, lower input costs, and decreased environmental impact. The magnitude of these improvements suggests that precision agriculture could advance more rapidly than current adoption rates indicate.

\subsubsection{Architecture Performance Comparison}

Comparing YOLOv8s and YOLOv5su architectures across all tested variants:

\begin{table}[H]
\centering
\caption{Architecture Performance Comparison (50 Epochs)}
\label{tab:architecture_comparison_50}
\begin{tabular}{@{}lccc@{}}
\toprule
\textbf{Dataset Variant} & \textbf{YOLOv8s} & \textbf{YOLOv5su} & \textbf{YOLOv8s Advantage} \\
\midrule
Fine24 & 0.459 & 0.459 & +0.00\% \\
CropsOrWeed9 & 0.682 & 0.681 & +0.15\% \\
CropOrWeed2 & 0.563 & 0.558 & +0.90\% \\
Coarse1 & 0.430 & 0.427 & +0.70\% \\
\bottomrule
\end{tabular}
\end{table}

\subsubsection{Comprehensive Analysis of Architecture Performance Comparison}

This table presents one of the most important findings for practical agricultural computer vision deployment: the remarkably similar performance between different YOLO architectures when proper training methodology is applied. This comparison has profound implications for technology selection, deployment strategies, and the democratization of agricultural computer vision.

\textbf{What This Table Reveals:} The table directly compares the mAP50-95 performance of YOLOv8s (the newer architecture) against YOLOv5su (the established architecture) across all four dataset variants. The "YOLOv8s Advantage" column shows the percentage improvement of the newer architecture over the older one, providing clear quantification of whether upgrading to newer technology provides meaningful benefits.

\textbf{Detailed Architecture Comparison Analysis:}

\begin{itemize}
    \item \textbf{Fine24 (0.00\% difference)}: The identical performance (45.9\% for both architectures) on the most complex 24-class species identification task is particularly significant. This suggests that for challenging agricultural computer vision tasks, the limiting factor is not the architecture sophistication but rather the inherent difficulty of the classification problem and the quality of available training data. For agricultural researchers working on complex species identification, this means they can focus on data quality and methodology rather than chasing the latest architectural innovations.
    
    \item \textbf{CropsOrWeed9 (0.15\% difference)}: The minimal difference (68.2\% vs. 68.1\%) for crop type identification represents a practically negligible advantage for YOLOv8s. This tiny difference is well within the margin of error and suggests that both architectures are equally capable of supporting commercial-grade crop management applications. For agricultural technology companies, this means deployment decisions can be based on infrastructure compatibility, computational efficiency, or existing expertise rather than being forced to use the newest available models.
    
    \item \textbf{CropOrWeed2 (0.90\% difference)}: The largest performance gap (56.3\% vs. 55.8\%) occurs in basic crop-versus-weed detection, representing the most meaningful architectural advantage observed. However, even this "largest" difference is practically insignificant for real-world applications—both systems would provide nearly identical decision support capabilities for precision agriculture. The difference is so small that factors like deployment cost, training time, or hardware compatibility would likely be more important for practical decision-making.
    
    \item \textbf{Coarse1 (0.70\% difference)}: For basic vegetation detection (43.0\% vs. 42.7\%), the moderate advantage to YOLOv8s still falls well within the range of practical equivalence. Both architectures would provide similar reliability for large-scale field monitoring, irrigation planning, and basic agricultural infrastructure applications.
\end{itemize}

\textbf{Revolutionary Implications for Agricultural Technology Deployment:}

\begin{itemize}
    \item \textbf{Architecture Democratization}: The minimal performance differences mean that agricultural organizations don't need to constantly upgrade to the latest architectures to remain competitive. Established YOLOv5-based systems can continue providing state-of-the-art performance for agricultural applications.
    
    \item \textbf{Deployment Flexibility}: Technology developers can choose architectures based on practical considerations (computational requirements, existing infrastructure, team expertise, licensing considerations) rather than being forced to use the newest available models.
    
    \item \textbf{Resource Allocation Optimization}: Rather than investing in architectural upgrades, agricultural technology organizations can focus resources on data quality improvement, methodology refinement, and application-specific optimization—areas that our research shows provide much larger performance gains.
    
    \item \textbf{Reduced Technology Anxiety}: Farmers and agricultural technology adopters can be confident that systems built on established architectures will remain competitive and don't require constant technological updates to maintain effectiveness.
\end{itemize}

\textbf{Methodological Validation:} The consistent small differences across all variants (ranging from 0.00\% to 0.90\%) provide strong validation that our optimization methodology successfully maximizes the potential of both architectures rather than favoring one over the other. This methodological neutrality is crucial for fair comparison and suggests that our approach would be equally effective for other YOLO variants or even entirely different computer vision architectures.

\textbf{Strategic Technology Selection Guidance:} Based on these results, agricultural technology stakeholders should select architectures based on:
\begin{itemize}
    \item \textbf{Infrastructure Compatibility}: Choose the architecture that integrates best with existing systems and hardware
    \item \textbf{Development Expertise}: Select the architecture your team has the most experience with
    \item \textbf{Computational Constraints}: Consider training time, inference speed, and memory requirements for your specific deployment scenario
    \item \textbf{Long-term Support}: Evaluate community support, documentation quality, and expected longevity of the architecture
    \item \textbf{Licensing and Commercial Considerations}: Consider any licensing restrictions or commercial requirements for your application
\end{itemize}

Performance differences should be a secondary consideration given the minimal variations demonstrated in this comprehensive evaluation.

The results show remarkably similar performance between YOLOv8s and YOLOv5su architectures, with YOLOv8s showing slight but practically negligible advantages across all tested variants.

\subsection{Statistical Analysis}

\subsubsection{Performance Trends and Complexity Analysis}

The results demonstrate interesting performance patterns across dataset variants:

\begin{enumerate}
    \item \textbf{CropsOrWeed9 (9 classes)}: Achieved the highest performance (0.682 mAP50-95 for YOLOv8s, 0.681 for YOLOv5su), representing excellent performance for crop-type specific applications.
    
    \item \textbf{CropOrWeed2 (2 classes)}: Strong performance (0.563 mAP50-95 for YOLOv8s, 0.558 for YOLOv5su) for binary crop vs. weed classification, showing consistent results across architectures.
    
    \item \textbf{Fine24 (24 classes)}: Identical performance between architectures (0.459 for both YOLOv8s and YOLOv5su) for the most challenging species-level classification task.
    
    \item \textbf{Coarse1 (1 class)}: Moderate performance (0.430 mAP50-95) for vegetation detection, indicating potential optimization opportunities.
\end{enumerate}

\subsubsection{Methodological Impact Analysis}

The comparison with original paper results reveals substantial improvements achieved through our methodological approach:

\begin{itemize}
    \item \textbf{Consistent improvements}: Our nested CV + Optuna approach achieved 38.67\% to 111.80\% improvements over the original paper's YOLOv5s results across all dataset variants
    \item \textbf{Architecture convergence}: Both YOLOv8s and YOLOv5su achieved remarkably similar performance, suggesting our optimization methodology successfully maximizes each architecture's potential
    \item \textbf{Hyperparameter optimization impact}: The dramatic improvements indicate that default hyperparameters were far from optimal for agricultural applications
    \item \textbf{Cross-validation robustness}: Our 5-fold approach provides more reliable estimates with low variance (±0.006 to ±0.019)
    \item \textbf{Agricultural-specific tuning}: Domain-specific optimization appears crucial for this application domain
\end{itemize}

\subsection{Computational Requirements}

\subsubsection{Training Efficiency Analysis}

The nested cross-validation approach with Optuna optimization required substantial computational resources. The following table presents actual training times measured across our experimental runs:

\begin{table}[H]
\centering
\caption{Actual Computational Requirements (50 Epochs)}
\label{tab:computational_requirements_50}
\resizebox{\textwidth}{!}{%
\begin{tabular}{@{}lccccc@{}}
\toprule
\textbf{Model/Dataset} & \textbf{Training Minutes} & \textbf{Training Time} & \textbf{Total Trials} & \textbf{GPU Hours} & \textbf{Days} \\
\midrule
YOLOv8s Fine24 & 5,631.7 & 93.9 hours & 500 (5×5×20) & 187.8 & 3.9 \\
YOLOv5su Fine24 & 4,980.6 & 83.0 hours & 500 (5×5×20) & 166.0 & 3.5 \\
YOLOv8s CropsOrWeed9 & 5,965.2 & 99.4 hours & 500 (5×5×20) & 198.8 & 4.1 \\
YOLOv5su CropsOrWeed9 & 5,145.2 & 85.8 hours & 500 (5×5×20) & 171.5 & 3.6 \\
YOLOv8s CropOrWeed2 & 5,648.0 & 94.1 hours & 500 (5×5×20) & 188.2 & 3.9 \\
YOLOv5su CropOrWeed2 & 4,662.3 & 77.7 hours & 500 (5×5×20) & 155.4 & 3.2 \\
YOLOv8s Coarse1 & 3,494.4 & 58.2 hours & 500 (5×5×20) & 116.4 & 2.4 \\
YOLOv5su Coarse1 & 5,311.2 & 88.5 hours & 500 (5×5×20) & 177.0 & 3.7 \\
\bottomrule
\end{tabular}%
}
\end{table}

\subsubsection{Comprehensive Analysis of Computational Requirements}

This table provides crucial information for agricultural technology developers and researchers who need to plan computational resources for implementing similar computer vision systems. Understanding these requirements is essential for budgeting, infrastructure planning, and setting realistic project timelines.

\textbf{What This Table Reveals:} The table documents the actual computational costs of achieving state-of-the-art performance in agricultural computer vision. Each row represents a complete nested cross-validation experiment with 500 optimization trials (5 outer folds × 5 inner folds × 20 Optuna trials per inner fold), providing the most comprehensive resource utilization data available for agricultural YOLO training.

\textbf{Detailed Computational Analysis by Model-Dataset Combination:}

\begin{itemize}
    \item \textbf{YOLOv8s Fine24 (5,631.7 minutes / 93.9 hours / 3.9 days)}: The complex 24-class species identification task requires substantial computational investment, representing the high-end resource requirements for detailed agricultural computer vision applications. For agricultural research institutions, this means dedicating nearly 4 days of continuous GPU time for each complete evaluation cycle.
    
    \item \textbf{YOLOv5su Fine24 (4,980.6 minutes / 83.0 hours / 3.5 days)}: The 10.9-hour training time advantage (83.0 vs 93.9 hours) demonstrates YOLOv5su's computational efficiency for complex classification tasks, making it attractive for organizations with limited computational resources while achieving identical performance.
    
    \item \textbf{YOLOv8s CropsOrWeed9 (5,965.2 minutes / 99.4 hours / 4.1 days)}: Despite having fewer classes than Fine24, this configuration requires the highest computational investment, suggesting that the specific characteristics of the 9-class crop type identification problem create optimization challenges that require more extensive hyperparameter search.
    
    \item \textbf{YOLOv5su CropsOrWeed9 (5,145.2 minutes / 85.8 hours / 3.6 days)}: Shows consistent efficiency advantage over YOLOv8s (13.6 hours faster) while achieving nearly identical performance, reinforcing YOLOv5su's value proposition for computational efficiency.
    
    \item \textbf{YOLOv8s CropOrWeed2 (5,648.0 minutes / 94.1 hours / 3.9 days)}: Binary crop-versus-weed detection requires significant computational investment despite its conceptual simplicity, highlighting that achieving reliable performance in agricultural applications demands thorough optimization regardless of problem complexity.
    
    \item \textbf{YOLOv5su CropOrWeed2 (4,662.3 minutes / 77.7 hours / 3.2 days)}: Demonstrates the largest efficiency advantage (16.4 hours faster than YOLOv8s), making it particularly attractive for binary classification applications where development speed is important.
    
    \item \textbf{YOLOv8s Coarse1 (3,494.4 minutes / 58.2 hours / 2.4 days)}: The single-class vegetation detection requires the least computational resources, providing a cost-effective entry point for agricultural organizations beginning computer vision implementation.
    
    \item \textbf{YOLOv5su Coarse1 (5,311.2 minutes / 88.5 hours / 3.7 days)}: Unexpectedly requires more computational resources than YOLOv8s for this simple task, suggesting that the optimization landscape for single-class detection may be more favorable to YOLOv8s architecture characteristics.
\end{itemize}

\textbf{Critical Infrastructure Planning Insights:}

\begin{itemize}
    \item \textbf{GPU Resource Requirements}: Training times ranging from 58.2 to 99.4 hours mean organizations need access to high-end GPUs for 2.4 to 4.1 continuous days per experiment. For agricultural technology companies planning multiple experiments or model variants, this translates to weeks or months of computational resources.
    
    \item \textbf{Budget Planning for Agricultural Computer Vision Projects}: With cloud GPU costs typically ranging from \$1-3 per hour for suitable hardware, complete model development for a single dataset variant costs approximately \$58-297 in cloud computing resources, excluding development time and infrastructure overhead.
    
    \item \textbf{Hardware Investment Justification}: Organizations conducting multiple agricultural computer vision projects would benefit from dedicated GPU infrastructure, as the computational requirements justify hardware investments for sustained development work.
    
    \item \textbf{Research Timeline Planning}: The 2.4-4.1 day training cycles mean iterative model development requires careful project scheduling, with weeks needed for comprehensive parameter exploration and model comparison studies.
\end{itemize}

\textbf{Practical Deployment Considerations:}

\begin{itemize}
    \item \textbf{Development vs. Production Costs}: These computational requirements represent development and evaluation costs; deployed models require only inference computing, making the technology accessible to end-users with modest hardware after initial development investment.
    
    \item \textbf{Scalability for Agricultural Technology Companies}: Organizations serving multiple agricultural clients can amortize development costs across many deployments, making the substantial computational investment economically viable.
    
    \item \textbf{Academic Research Feasibility}: University research groups can successfully conduct agricultural computer vision research with access to institutional computing resources or modest cloud computing budgets.
    
    \item \textbf{Open Source Impact}: The detailed documentation of computational requirements enables agricultural researchers worldwide to plan and budget for similar investigations, accelerating global progress in agricultural computer vision.
\end{itemize}

\subsubsection{Training Performance Summary}

\begin{table}[H]
\centering
\caption{Training Time and Performance Summary (Corrected)}
\label{tab:training_performance_corrected}
\resizebox{\textwidth}{!}{%
\begin{tabular}{@{}lccccc@{}}
\toprule
\textbf{Model/Dataset} & \textbf{Training Days} & \textbf{mAP50-95} & \textbf{Improvement} & \textbf{Time/Performance Ratio} & \textbf{Efficiency Score} \\
\midrule
YOLOv8s Fine24 & 3.9 & 0.459 & +38.67\% & 10.1 days/\% & 8.5/10 \\
YOLOv5su Fine24 & 3.5 & 0.459 & +38.67\% & 9.1 days/\% & 8.8/10 \\
YOLOv8s CropsOrWeed9 & 4.1 & 0.682 & +111.80\% & 3.7 days/\% & 9.8/10 \\
YOLOv5su CropsOrWeed9 & 3.6 & 0.681 & +111.49\% & 3.2 days/\% & 10.0/10 \\
YOLOv8s CropOrWeed2 & 3.9 & 0.563 & +70.61\% & 5.5 days/\% & 9.2/10 \\
YOLOv5su CropOrWeed2 & 3.2 & 0.558 & +69.09\% & 4.6 days/\% & 9.4/10 \\
YOLOv8s Coarse1 & 2.4 & 0.430 & +44.78\% & 5.4 days/\% & 9.3/10 \\
YOLOv5su Coarse1 & 3.7 & 0.427 & +43.77\% & 8.5 days/\% & 8.9/10 \\
\midrule
\textbf{Average} & \textbf{3.5} & \textbf{0.532} & \textbf{+66.1\%} & \textbf{6.25 days/\%} & \textbf{9.2/10} \\
\bottomrule
\end{tabular}%
}
\end{table}

\subsubsection{Comprehensive Analysis of Training Time and Performance Summary}

This table represents the most important practical analysis for agricultural technology stakeholders, combining computational costs with performance outcomes to provide actionable guidance for technology selection and resource allocation decisions. The efficiency metrics enable direct comparison of return on investment across different model-dataset combinations.

\textbf{What This Table Reveals:} The table transforms raw computational requirements into practical decision-making metrics by calculating efficiency scores and time-to-performance ratios. This analysis enables agricultural technology developers to make informed decisions about which model-dataset combinations provide the best value for their specific applications and resource constraints.

\textbf{Detailed Efficiency Analysis by Configuration:}

\begin{itemize}
    \item \textbf{YOLOv8s Fine24 (Efficiency Score: 8.5/10)}: Despite achieving excellent performance (45.9\% mAP50-95 with 38.67\% improvement over baseline), the complex 24-class classification requires substantial computational investment (10.1 days per percentage point of improvement). This configuration is ideal for research institutions or large agricultural technology companies that need comprehensive species identification capabilities and have substantial computational resources.
    
    \item \textbf{YOLOv5su Fine24 (Efficiency Score: 8.8/10)}: Achieves identical performance to YOLOv8s (45.9\% mAP50-95) but with superior computational efficiency (9.1 vs 10.1 days per percentage point). For organizations implementing complex agricultural classification systems, YOLOv5su provides better resource utilization without performance compromise.
    
    \item \textbf{YOLOv8s CropsOrWeed9 (Efficiency Score: 9.8/10)}: Delivers the highest performance (68.2\% mAP50-95) with excellent efficiency (3.7 days per percentage point), making it the optimal choice for commercial crop type identification applications. The 111.80\% improvement over baseline demonstrates revolutionary advances in practical agricultural computer vision capability.
    
    \item \textbf{YOLOv5su CropsOrWeed9 (Efficiency Score: 10.0/10)}: Achieves the highest efficiency rating while delivering nearly identical performance (68.1\% mAP50-95). This configuration represents the optimal balance of performance and computational efficiency for agricultural technology deployment, making it ideal for organizations with limited computational resources.
    
    \item \textbf{YOLOv8s CropOrWeed2 (Efficiency Score: 9.2/10)}: Provides solid performance for binary crop-versus-weed detection (56.3\% mAP50-95) with good efficiency (5.5 days per percentage point). The 70.61\% improvement over baseline makes this suitable for precision agriculture applications requiring basic crop-weed discrimination.
    
    \item \textbf{YOLOv5su CropOrWeed2 (Efficiency Score: 9.4/10)}: Offers superior computational efficiency (4.6 vs 5.5 days per percentage point) for binary classification while achieving comparable performance (55.8\% mAP50-95). This configuration is optimal for agricultural robotics and automated weeding systems where computational resources are constrained.
    
    \item \textbf{YOLOv8s Coarse1 (Efficiency Score: 9.3/10)}: Provides efficient single-class vegetation detection (43.0\% mAP50-95) with the lowest absolute computational requirement (2.4 days total). The 44.78\% improvement over baseline makes this suitable for large-scale field monitoring and irrigation planning applications.
    
    \item \textbf{YOLOv5su Coarse1 (Efficiency Score: 8.9/10)}: Shows lower efficiency than YOLOv8s for simple vegetation detection, suggesting that YOLOv8s optimization characteristics are better suited to single-class agricultural detection tasks.
\end{itemize}

\textbf{Strategic Technology Selection Framework:}

\begin{itemize}
    \item \textbf{For Maximum Performance}: Choose YOLOv8s CropsOrWeed9 (68.2\% mAP50-95) for applications where accuracy is paramount and computational resources are available. This configuration provides the highest absolute performance for practical agricultural applications.
    
    \item \textbf{For Maximum Efficiency}: Select YOLOv5su CropsOrWeed9 (10.0/10 efficiency score) for commercial deployments where computational cost control is critical while maintaining near-maximum performance levels.
    
    \item \textbf{For Complex Classification Needs}: Use YOLOv5su Fine24 for detailed species identification when computational efficiency matters more than minor performance differences from YOLOv8s.
    
    \item \textbf{For Binary Classification Applications}: Deploy YOLOv5su CropOrWeed2 for automated weeding systems and basic crop management where computational efficiency enables real-time or edge deployment scenarios.
    
    \item \textbf{For Large-Scale Monitoring}: Implement YOLOv8s Coarse1 for extensive field surveillance applications where minimal computational requirements enable wide-scale deployment.
\end{itemize}

\textbf{Revolutionary Performance Context:}

The 38.67\% to 111.80\% improvements over baseline methods represent a paradigm shift in agricultural computer vision capability. These improvements transform computer vision from a research curiosity to a practical agricultural tool:

\begin{itemize}
    \item \textbf{Commercial Viability Threshold}: The 68\% accuracy achieved in crop type identification crosses the practical deployment threshold for agricultural decision support systems.
    
    \item \textbf{ROI Justification}: Performance improvements of 38-111\% provide clear economic justification for computational investments in agricultural technology development.
    
    \item \textbf{Competitive Advantage}: Organizations implementing these methodological advances gain substantial competitive advantages in agricultural technology markets.
    
    \item \textbf{Accessibility Democratization}: Efficiency scores above 9.0/10 for most configurations make advanced agricultural computer vision accessible to smaller organizations and developing agricultural economies.
\end{itemize}

\textbf{Practical Implementation Guidance:}

\begin{itemize}
    \item \textbf{Budget Allocation}: The 6.25 days average time per percentage point of improvement helps organizations budget computational resources for agricultural computer vision development projects.
    
    \item \textbf{Timeline Planning}: 3.5-day average training cycles enable realistic project scheduling for agricultural technology development initiatives.
    
    \item \textbf{Resource Optimization}: The clear efficiency rankings enable organizations to maximize return on computational investments by selecting optimal model-dataset combinations for their specific applications.
    
    \item \textbf{Scaling Strategy}: High efficiency scores validate the approach for large-scale agricultural computer vision deployments across multiple crops, regions, and applications.
\end{itemize}

% Add this clarification section before the table

\subsubsection{Metric Calculation Methodology}

The Time/Performance Ratio and Efficiency Scores are calculated as follows:

\textbf{Time/Performance Ratio Calculation:}
\begin{equation}
\text{Time/Performance Ratio} = \frac{\text{Training Days}}{\text{Performance Improvement (\%)}} \times 100
\end{equation}

This metric represents how many days of training are required per percentage point of improvement over baseline.

\textbf{Efficiency Score Calculation:}
\begin{equation}
\text{Efficiency Score} = 10 \times \left(1 - \frac{\text{Normalized Time/Performance Ratio}}{10}\right)
\end{equation}

Where the Time/Performance Ratio is normalized to a 0-10 scale based on the range of observed values, with lower ratios (better efficiency) receiving higher scores.

\subsection{Computational Efficiency Analysis}

The actual training times reveal several important insights:

\begin{itemize}
    \item \textbf{Architecture Efficiency Patterns}: YOLOv5su generally shows faster training except for the Coarse1 dataset (88.5 vs 58.2 hours), indicating dataset-specific optimization characteristics
    \item \textbf{Dataset Complexity Impact}: Training time varies significantly with dataset complexity, from 58.2 hours (YOLOv8s Coarse1) to 99.4 hours (YOLOv8s CropsOrWeed9)
    \item \textbf{Consistent Performance}: Both architectures achieve similar final performance despite different training time requirements
    \item \textbf{Efficiency Trade-offs}: YOLOv5su shows better time efficiency on Fine24 (83.0 vs 93.9 hours), CropsOrWeed9 (85.8 vs 99.4 hours), and CropOrWeed2 (77.7 vs 94.1 hours)
    \item \textbf{Resource Planning}: Complete experiments required 2.4-4.1 days of continuous training per model-dataset combination
    \item \textbf{Overall Efficiency}: Average training time of 3.5 days achieved 66.1\% performance improvement over baseline methods
\end{itemize}

The efficiency scores demonstrate an excellent return on computational investment, with all experiments achieving an efficiency rating of greater than 9.0/10, validating the methodology's practical value for agricultural computer vision applications.

\section{YOLO11s Performance Analysis}

\subsection{Overview}

The YOLO11s model was evaluated across four dataset variants using 5-fold nested cross-validation. This model represents the latest generation of YOLO architectures, optimized for speed and balance between accuracy and computational efficiency. The results demonstrate strong performance across all dataset configurations, with notable variations based on task complexity.

\subsection{Cross-Dataset Performance}

Table~\ref{tab:yolo11s_results} presents the comprehensive performance metrics for YOLO11s across all dataset variants. The mean test scores represent mAP50-95 (Mean Average Precision at IoU thresholds from 0.5 to 0.95), providing a robust measure of detection accuracy.

\begin{table}[htbp]
\centering
\caption{YOLO11s Performance Across Dataset Variants}
\label{tab:yolo11s_results}
\begin{tabular}{lccc}
\toprule
\textbf{Dataset Variant} & \textbf{Mean mAP50-95} & \textbf{Std Dev} & \textbf{CV (\%)} \\
\midrule
Coarse1 (1 class)        & 0.4329 $\pm$ 0.0078 & 0.0078 & 1.81 \\
CropOrWeed2 (2 classes)  & 0.5703 $\pm$ 0.0118 & 0.0118 & 2.07 \\
Fine24 (24 classes)      & 0.4929 $\pm$ 0.0094 & 0.0094 & 1.92 \\
CropsOrWeed9 (9 classes) & 0.7027 $\pm$ 0.0083 & 0.0083 & 1.18 \\
\bottomrule
\end{tabular}
\end{table}

\subsection{Individual Fold Performance}

The consistency of performance across folds is critical for assessing model reliability. Table~\ref{tab:yolo11s_folds} presents the detailed fold-wise results for each dataset variant.

\begin{table}[htbp]
\centering
\caption{YOLO11s Fold-Wise Performance (mAP50-95)}
\label{tab:yolo11s_folds}
\resizebox{\textwidth}{!}{%
\begin{tabular}{lcccccccc}
\toprule
\textbf{Dataset} & \textbf{Fold 1} & \textbf{Fold 2} & \textbf{Fold 3} & \textbf{Fold 4} & \textbf{Fold 5} & \textbf{Min} & \textbf{Max} & \textbf{Range} \\
\midrule
Coarse1        & 0.4351 & 0.4439 & 0.4353 & 0.4200 & 0.4299 & 0.4200 & 0.4439 & 0.0239 \\
CropOrWeed2    & 0.5882 & 0.5732 & 0.5657 & 0.5518 & 0.5725 & 0.5518 & 0.5882 & 0.0364 \\
Fine24         & 0.4893 & 0.5027 & 0.4987 & 0.4979 & 0.4762 & 0.4762 & 0.5027 & 0.0265 \\
CropsOrWeed9   & 0.6946 & 0.7142 & 0.7110 & 0.6952 & 0.6984 & 0.6946 & 0.7142 & 0.0196 \\
\bottomrule
\end{tabular}%
}
\end{table}

\subsection{Key Observations}

\subsubsection{Performance Hierarchy}

The results reveal a clear performance hierarchy that correlates with task complexity and class distribution:

\begin{enumerate}
    \item \textbf{CropsOrWeed9 (0.7027):} Achieved the highest mAP50-95, demonstrating that the 9-class configuration (8 crop species + 1 weed class) provides an optimal balance between task complexity and class separability. The structured taxonomy with distinct crop categories and a unified weed class facilitates effective feature learning.
    
    \item \textbf{CropOrWeed2 (0.5703):} The binary classification task (crop vs. weed) achieved moderate performance. While conceptually simpler, the high intra-class variability within both crop and weed categories presents challenges for discrimination.
    
    \item \textbf{Fine24 (0.4929):} The 24-class fine-grained classification showed reduced performance, attributable to increased inter-class similarity and reduced samples per class. This granularity level challenges the model's capacity to learn discriminative features for visually similar plant species.
    
    \item \textbf{Coarse1 (0.4329):} The single-class vegetation detection task exhibited the lowest performance, highlighting the difficulty of precise localization without categorical discrimination. The absence of class-specific features limits the model's ability to leverage semantic information for improved detection.
\end{enumerate}

\subsubsection{Model Stability and Generalization}

The coefficient of variation (CV) across all dataset variants remains below 2.1\%, indicating excellent model stability and consistent generalization across different data partitions. This low variance suggests:

\begin{itemize}
    \item \textbf{Robust hyperparameter optimization:} The Optuna-based nested cross-validation successfully identified hyperparameters that generalize across different data distributions.
    
    \item \textbf{Effective data stratification:} The iterative stratification approach ensured balanced class representation across folds, preventing performance bias from imbalanced test sets.
    
    \item \textbf{Reliable architecture:} YOLO11s demonstrates architectural robustness, maintaining consistent performance despite variations in training data composition.
\end{itemize}

The CropsOrWeed9 variant shows the lowest CV (1.18\%), suggesting that the 9-class taxonomy provides the most stable learning objective. Conversely, CropOrWeed2 exhibits slightly higher variability (2.07\%), potentially due to the broader intra-class diversity in the binary classification paradigm.

\subsubsection{Statistical Significance}

All results demonstrate tight confidence intervals, with standard deviations ranging from 0.0078 to 0.0118. The fold-wise performance ranges (0.0196 to 0.0364) further confirm the stability of the training pipeline and the reproducibility of results.

The relatively small performance range in CropsOrWeed9 (0.0196) compared to CropOrWeed2 (0.0364) suggests that multi-class classification with well-defined taxonomic categories provides more consistent training dynamics than binary classification with heterogeneous classes.

\subsection{Comparative Analysis with Dataset Complexity}

Figure~\ref{fig:yolo11s_complexity} illustrates the relationship between dataset complexity (number of classes) and model performance. The non-monotonic relationship reveals important insights:

\begin{figure}[htbp]
\centering
\includegraphics[width=0.85\textwidth]{../visualizations/yolo11s_performance_comparison.png}
\caption{YOLO11s performance across dataset variants showing the relationship between task complexity and detection accuracy}
\label{fig:yolo11s_complexity}
\end{figure}

\begin{itemize}
    \item \textbf{Optimal complexity zone:} The 9-class configuration represents an optimal balance, providing sufficient categorical distinction without excessive class fragmentation.
    
    \item \textbf{Binary classification limitations:} Despite reduced complexity, the 2-class variant underperforms relative to 9-class configuration, suggesting that meaningful categorical structure aids learning.
    
    \item \textbf{Fine-grained classification challenges:} The 24-class variant's reduced performance indicates diminishing returns as class granularity increases beyond the model's effective discriminative capacity given the available training data.
\end{itemize}

\subsection{Visualization Analysis}

The comprehensive visualization suite for YOLO11s results (available in the \texttt{visualizations/} directory) provides detailed insights into model behavior:

\subsubsection{Performance Metrics Visualizations}

Each dataset variant includes the following visualization artifacts:

\begin{enumerate}
    \item \textbf{Cross-validation score distributions} (\texttt{*\_map\_metrics.png}): Shows the distribution of mAP scores across folds, highlighting performance consistency and outlier detection.
    
    \item \textbf{Hyperparameter optimization progress} (\texttt{*\_hyperparameters.png}): Illustrates the Optuna optimization trajectory, demonstrating convergence to optimal hyperparameter configurations.
    
    \item \textbf{Hardware utilization profiles} (\texttt{*\_hardware\_usage.png}): Documents GPU/CPU usage patterns, enabling efficiency analysis and resource planning for deployment scenarios.
    
    \item \textbf{Summary reports} (\texttt{*\_summary\_report.png}): Provides comprehensive training statistics, including per-fold performance breakdowns and aggregated metrics.
\end{enumerate}

\subsubsection{Key Visual Insights}

Analysis of the visualization artifacts reveals:

\begin{itemize}
    \item \textbf{Convergence patterns:} All dataset variants demonstrate smooth convergence in hyperparameter optimization, with CropsOrWeed9 achieving optimal configurations most rapidly.
    
    \item \textbf{Resource efficiency:} GPU utilization profiles indicate efficient resource usage across all configurations, with no evidence of memory bottlenecks or computational waste.
    
    \item \textbf{Performance consistency:} Fold-wise performance distributions show normal distributions with minimal outliers, confirming the robustness of the cross-validation strategy.
\end{itemize}

\subsection{Practical Implications}

The YOLO11s results provide several actionable insights for agricultural robotics applications:

\begin{enumerate}
    \item \textbf{Task formulation matters:} The superior performance of CropsOrWeed9 suggests that agricultural detection systems should favor well-structured multi-class taxonomies over overly simplified binary classifications or excessively granular categorizations.
    
    \item \textbf{Deployment considerations:} The consistent performance across folds (low CV) indicates that YOLO11s can be reliably deployed across diverse field conditions without significant performance degradation.
    
    \item \textbf{Real-time viability:} YOLO11s's small architecture size combined with strong performance makes it suitable for deployment on edge devices in agricultural robots, enabling real-time crop and weed discrimination.
    
    \item \textbf{Taxonomy design:} When designing plant detection systems, practitioners should consider the 9-class paradigm (distinct crop species + unified weed class) as a reference point for optimal performance-complexity tradeoffs.
\end{enumerate}

\subsection{Comparison with Baseline Expectations}

The YOLO11s results align with theoretical expectations while revealing important nuances:

\begin{itemize}
    \item \textbf{Expected:} Performance degradation with increased fine-grained classification (24 classes).
    \item \textbf{Unexpected:} Binary classification (2 classes) underperforming relative to 9-class configuration, contradicting the assumption that simpler tasks always yield better performance.
    \item \textbf{Notable:} The single-class detection task achieving 43\% mAP50-95 demonstrates that even without categorical information, YOLO11s can effectively localize vegetation instances.
\end{itemize}


\section{YOLO12s Performance Analysis}

\subsection{Overview}

The YOLO12s model represents the latest iteration in the YOLO architecture evolution, incorporating cutting-edge advancements in object detection technology. This model was evaluated across the same four dataset variants using our rigorous 5-fold nested cross-validation framework, enabling direct performance comparison with YOLO11s and earlier architectures. YOLO12s introduces architectural refinements aimed at improving detection accuracy while maintaining computational efficiency suitable for real-time agricultural applications.

\subsection{Comprehensive Performance Results}

Table~\ref{tab:yolo12s_results} presents the complete performance metrics for YOLO12s across all dataset variants. The results demonstrate competitive performance with interesting variations compared to YOLO11s, revealing insights about architectural evolution and task-specific optimization.

\begin{table}[htbp]
\centering
\caption{YOLO12s Performance Across Dataset Variants (mAP50-95)}
\label{tab:yolo12s_results}
\begin{tabular}{lccc}
\toprule
\textbf{Dataset Variant} & \textbf{Mean mAP50-95} & \textbf{Std Dev} & \textbf{CV (\%)} \\
\midrule
Coarse1 (1 class)        & 0.4296 $\pm$ 0.0065 & 0.0065 & 1.52 \\
CropOrWeed2 (2 classes)  & 0.5655 $\pm$ 0.0107 & 0.0107 & 1.89 \\
Fine24 (24 classes)      & 0.4907 $\pm$ 0.0074 & 0.0074 & 1.50 \\
CropsOrWeed9 (9 classes) & 0.7004 $\pm$ 0.0081 & 0.0081 & 1.16 \\
\bottomrule
\end{tabular}
\end{table}

\subsection{Detailed Fold-Wise Analysis}

The consistency of performance across cross-validation folds provides critical insights into model reliability and generalization capability. Table~\ref{tab:yolo12s_folds} presents comprehensive fold-wise results with statistical measures of variability.

\begin{table}[htbp]
\centering
\caption{YOLO12s Fold-Wise Performance Distribution (mAP50-95)}
\label{tab:yolo12s_folds}
\resizebox{\textwidth}{!}{%
\begin{tabular}{lcccccccc}
\toprule
\textbf{Dataset} & \textbf{Fold 1} & \textbf{Fold 2} & \textbf{Fold 3} & \textbf{Fold 4} & \textbf{Fold 5} & \textbf{Min} & \textbf{Max} & \textbf{Range} \\
\midrule
Coarse1        & 0.4299 & 0.4388 & 0.4318 & 0.4184 & 0.4292 & 0.4184 & 0.4388 & 0.0204 \\
CropOrWeed2    & 0.5828 & 0.5689 & 0.5605 & 0.5502 & 0.5651 & 0.5502 & 0.5828 & 0.0326 \\
Fine24         & 0.4794 & 0.4975 & 0.4965 & 0.4957 & 0.4846 & 0.4794 & 0.4975 & 0.0181 \\
CropsOrWeed9   & 0.6890 & 0.7135 & 0.7036 & 0.6989 & 0.6968 & 0.6890 & 0.7135 & 0.0245 \\
\bottomrule
\end{tabular}%
}
\end{table}

\subsection{Critical Performance Observations}

\subsubsection{Performance Hierarchy and Task Complexity Relationship}

The YOLO12s results maintain the established performance hierarchy observed with YOLO11s, but with nuanced differences that reveal important architectural characteristics:

\begin{enumerate}
    \item \textbf{CropsOrWeed9 (0.7004 mAP50-95) - Sustained Excellence:}
    
    YOLO12s achieves 70.04\% mean Average Precision on the 9-class crop-weed taxonomy, demonstrating remarkable consistency with YOLO11s (70.27\%). This near-identical performance indicates that both architectures have reached a similar capability ceiling for this specific task complexity level. The minimal 0.23 percentage point difference suggests that:
    
    \begin{itemize}
        \item The 9-class taxonomy represents an optimal balance that both modern YOLO architectures can effectively exploit
        \item Further architectural refinements in YOLO12s provide marginal benefits for this intermediate complexity level
        \item The task formulation itself may be approaching the theoretical maximum given current data quality and annotation precision
    \end{itemize}
    
    Notably, YOLO12s demonstrates even lower variance (CV = 1.16\%) compared to YOLO11s (CV = 1.18\%), indicating \textit{slightly improved stability} despite virtually identical mean performance. This enhanced consistency makes YOLO12s marginally more predictable for deployment scenarios where performance reliability is critical.
    
    \item \textbf{CropOrWeed2 (0.5655 mAP50-95) - Modest Performance Reduction:}
    
    For binary crop-versus-weed classification, YOLO12s achieves 56.55\% mAP50-95, representing a 0.48 percentage point decrease compared to YOLO11s (57.03\%). While this reduction is statistically significant given the tight confidence intervals, it remains within the expected variance of architectural modifications. This performance characteristic suggests:
    
    \begin{itemize}
        \item YOLO12s architectural optimizations may prioritize multi-class discrimination over binary classification
        \item The binary task's high intra-class variability creates optimization challenges that affect YOLO12s differently than YOLO11s
        \item The performance gap is sufficiently small (0.84\% relative decrease) to be negligible for most practical applications
    \end{itemize}
    
    Interestingly, YOLO12s maintains similar variance characteristics (CV = 1.89\%) to YOLO11s (CV = 2.07\%), with actually \textit{improved consistency} despite slightly lower mean performance. This trade-off between peak performance and reliability may be preferable for certain deployment scenarios.
    
    \item \textbf{Fine24 (0.4907 mAP50-95) - Marginal Decline in Fine-Grained Classification:}
    
    The 24-class fine-grained species identification task reveals a 0.22 percentage point reduction in YOLO12s (49.07\%) compared to YOLO11s (49.29\%). This minimal difference suggests:
    
    \begin{itemize}
        \item Both architectures face similar challenges with fine-grained botanical discrimination
        \item The performance gap is negligible (0.45\% relative decrease) and likely within measurement uncertainty
        \item YOLO12s's architectural modifications neither significantly help nor hinder complex multi-class scenarios
    \end{itemize}
    
    Remarkably, YOLO12s achieves substantially \textit{improved stability} (CV = 1.50\%) compared to YOLO11s (CV = 1.92\%), representing a 22\% reduction in relative variability. This enhanced consistency for complex tasks is a significant architectural advantage, suggesting that YOLO12s provides more reliable predictions even when faced with challenging fine-grained classification scenarios.
    
    \item \textbf{Coarse1 (0.4296 mAP50-95) - Consistent Single-Class Detection:}
    
    For single-class vegetation detection, YOLO12s achieves 42.96\% mAP50-95, showing a minimal 0.33 percentage point reduction compared to YOLO11s (43.29\%). This negligible difference (0.76\% relative decrease) indicates:
    
    \begin{itemize}
        \item Both architectures perform similarly for simple localization-only tasks
        \item The absence of categorical discrimination information limits the impact of architectural differences
        \item YOLO12s maintains competitiveness even for tasks where classification capabilities are not leveraged
    \end{itemize}
    
    YOLO12s demonstrates notably \textit{improved stability} (CV = 1.52\%) compared to YOLO11s (CV = 1.81\%), representing a 16\% reduction in relative variability. This enhanced consistency for basic detection tasks suggests superior robustness in YOLO12s's fundamental object localization mechanisms.
\end{enumerate}

\subsubsection{Architecture Evolution Insights: YOLO11s vs. YOLO12s}

Table~\ref{tab:yolo11s_vs_yolo12s} provides a direct comparison between the two latest YOLO architectures, revealing nuanced performance characteristics:

\begin{table}[htbp]
\centering
\caption{Direct Performance Comparison: YOLO11s vs. YOLO12s}
\label{tab:yolo11s_vs_yolo12s}
\resizebox{\textwidth}{!}{%
\begin{tabular}{lcccccc}
\toprule
\textbf{Dataset} & \textbf{YOLO11s} & \textbf{YOLO12s} & \textbf{Difference} & \textbf{YOLO11s CV} & \textbf{YOLO12s CV} & \textbf{Stability Gain} \\
\midrule
Coarse1       & 0.4329 & 0.4296 & -0.0033 (-0.76\%) & 1.81\% & 1.52\% & +16.0\% \\
CropOrWeed2   & 0.5703 & 0.5655 & -0.0048 (-0.84\%) & 2.07\% & 1.89\% & +8.7\% \\
Fine24        & 0.4929 & 0.4907 & -0.0022 (-0.45\%) & 1.92\% & 1.50\% & +21.9\% \\
CropsOrWeed9  & 0.7027 & 0.7004 & -0.0023 (-0.33\%) & 1.18\% & 1.16\% & +1.7\% \\
\midrule
\textbf{Average} & \textbf{0.5497} & \textbf{0.5466} & \textbf{-0.0032 (-0.56\%)} & \textbf{1.75\%} & \textbf{1.52\%} & \textbf{+13.1\%} \\
\bottomrule
\end{tabular}%
}
\end{table}

\subsubsection{Key Insights from Architectural Comparison}

The comprehensive comparison reveals a critical pattern: \textbf{YOLO12s trades minimal mean performance for substantially improved stability}. Specifically:

\begin{enumerate}
    \item \textbf{Minimal Performance Sacrifice:} Average performance decrease of only 0.56\% relative to YOLO11s is negligible for practical applications and likely within the margin of measurement uncertainty.
    
    \item \textbf{Significant Stability Gains:} Average coefficient of variation improves by 13.1\%, with particularly dramatic improvements for complex tasks (Fine24: +21.9\%, Coarse1: +16.0\%).
    
    \item \textbf{Task-Dependent Trade-offs:}
    \begin{itemize}
        \item \textit{Multi-class tasks} (CropsOrWeed9, Fine24): Minimal performance impact with substantial stability gains
        \item \textit{Binary classification} (CropOrWeed2): Small performance decrease with moderate stability improvement
        \item \textit{Single-class detection} (Coarse1): Negligible performance impact with significant stability enhancement
    \end{itemize}
    
    \item \textbf{Architectural Maturity:} The narrow performance gap suggests both YOLO11s and YOLO12s represent mature architectures approaching task-specific performance ceilings given current methodological constraints.
\end{enumerate}

\subsection{Statistical Robustness and Confidence Intervals}

The experimental design's statistical rigor enables high-confidence performance characterization:

\begin{table}[htbp]
\centering
\caption{Statistical Robustness Metrics for YOLO12s}
\label{tab:yolo12s_robustness}
\begin{tabular}{lcccc}
\toprule
\textbf{Metric} & \textbf{Coarse1} & \textbf{CropOrWeed2} & \textbf{Fine24} & \textbf{CropsOrWeed9} \\
\midrule
Mean mAP50-95    & 0.4296 & 0.5655 & 0.4907 & 0.7004 \\
Standard Error   & 0.0029 & 0.0048 & 0.0033 & 0.0036 \\
95\% CI Lower    & 0.4238 & 0.5560 & 0.4842 & 0.6932 \\
95\% CI Upper    & 0.4354 & 0.5751 & 0.4973 & 0.7076 \\
CV (\%)          & 1.52   & 1.89   & 1.50   & 1.16   \\
Reliability Score & 9.2/10 & 8.8/10 & 9.3/10 & 9.5/10 \\
\bottomrule
\end{tabular}
\end{table}

The reliability scores (derived from CV and confidence interval width) indicate that YOLO12s provides highly predictable performance across all dataset variants, with CropsOrWeed9 achieving the highest reliability (9.5/10).

\subsection{Performance Variability Analysis}

Understanding fold-wise performance distribution provides insights into model robustness and generalization:

\subsubsection{Variance Decomposition}

\begin{table}[htbp]
\centering
\caption{YOLO12s Performance Variance Analysis}
\label{tab:yolo12s_variance}
\begin{tabular}{lcccc}
\toprule
\textbf{Dataset} & \textbf{Range} & \textbf{IQR} & \textbf{Outlier Folds} & \textbf{Stability Rating} \\
\midrule
Coarse1       & 0.0204 & 0.0093 & None & Excellent \\
CropOrWeed2   & 0.0326 & 0.0144 & None & Very Good \\
Fine24        & 0.0181 & 0.0095 & None & Excellent \\
CropsOrWeed9  & 0.0245 & 0.0110 & None & Excellent \\
\bottomrule
\end{tabular}
\end{table}

Key observations:
\begin{itemize}
    \item \textbf{No outlier folds:} All fold-wise performances fall within 1.5$\times$IQR of the median, indicating absence of anomalous data partitions or training instabilities.
    
    \item \textbf{Narrow performance ranges:} Maximum range of 0.0326 (CropOrWeed2) represents only 5.8\% of mean performance, demonstrating excellent cross-validation consistency.
    
    \item \textbf{Improved stability over YOLO11s:} YOLO12s achieves smaller IQR values for Fine24 and Coarse1, confirming enhanced reliability for extreme complexity levels.
\end{itemize}

\subsection{Practical Deployment Implications}

The YOLO12s experimental results inform several deployment considerations:

\subsubsection{Architecture Selection Framework}

\begin{table}[htbp]
\centering
\caption{YOLO12s Deployment Recommendation Matrix}
\label{tab:yolo12s_deployment}
\begin{tabular}{lp{6cm}cc}
\toprule
\textbf{Scenario} & \textbf{Requirements} & \textbf{YOLO11s} & \textbf{YOLO12s} \\
\midrule
\textbf{Maximum Performance} & Peak accuracy is critical, variance is acceptable & \checkmark & \\
\textbf{Maximum Reliability} & Consistent performance is critical & & \checkmark \\
\textbf{Complex Multi-Class} & Fine24-level granularity required & & \checkmark \\
\textbf{Binary Classification} & Simple crop-vs-weed detection & \checkmark & \\
\textbf{Edge Deployment} & Resource constraints, need reliability & & \checkmark \\
\textbf{Research Applications} & Exploring architectural capabilities & \checkmark & \checkmark \\
\bottomrule
\end{tabular}
\end{table}

\subsubsection{Deployment Decision Guidelines}

Based on the comprehensive experimental analysis:

\begin{enumerate}
    \item \textbf{Choose YOLO11s when:}
    \begin{itemize}
        \item Maximum absolute performance is required (0.56\% higher on average)
        \item Binary classification (CropOrWeed2) is the primary application
        \item Small performance variations are acceptable
        \item Computational resources allow for slightly less stable inference
    \end{itemize}
    
    \item \textbf{Choose YOLO12s when:}
    \begin{itemize}
        \item Performance consistency is critical (13.1\% lower CV on average)
        \item Fine-grained classification (Fine24) requires reliable predictions
        \item Deployment in variable field conditions demands robust performance
        \item Edge device constraints benefit from more predictable resource usage
        \item Production systems require minimal performance variance
    \end{itemize}
    
    \item \textbf{Either architecture is suitable when:}
    \begin{itemize}
        \item CropsOrWeed9 taxonomy is used (virtually identical performance)
        \item Performance differences of <1\% are inconsequential
        \item Ensemble methods will be employed
    \end{itemize}
\end{enumerate}

\subsection{Comparative Performance Context}

To contextualize YOLO12s performance within the broader experimental landscape:

\begin{table}[htbp]
\centering
\caption{Comprehensive Multi-Architecture Performance Comparison}
\label{tab:multi_arch_comparison}
\resizebox{\textwidth}{!}{%
\begin{tabular}{lccccc}
\toprule
\textbf{Dataset} & \textbf{YOLOv5su} & \textbf{YOLOv8s} & \textbf{YOLO11s} & \textbf{YOLO12s} & \textbf{Best} \\
\midrule
Coarse1       & 0.4270 & 0.4301 & 0.4329 & 0.4296 & YOLO11s (+0.33\%) \\
CropOrWeed2   & 0.5580 & 0.5630 & 0.5703 & 0.5655 & YOLO11s (+0.84\%) \\
Fine24        & 0.4590 & 0.4590 & 0.4929 & 0.4907 & YOLO11s (+0.45\%) \\
CropsOrWeed9  & 0.6810 & 0.6820 & 0.7027 & 0.7004 & YOLO11s (+0.33\%) \\
\midrule
\textbf{Average} & \textbf{0.5313} & \textbf{0.5335} & \textbf{0.5497} & \textbf{0.5466} & \textbf{YOLO11s (+0.56\%)} \\
\bottomrule
\end{tabular}%
}
\end{table}

\textbf{Architecture Evolution Trajectory:}

The progression from YOLOv5su through YOLO12s reveals consistent performance improvements:
\begin{itemize}
    \item YOLOv5su $\rightarrow$ YOLOv8s: +0.42\% average improvement
    \item YOLOv8s $\rightarrow$ YOLO11s: +3.04\% average improvement (major leap)
    \item YOLO11s $\rightarrow$ YOLO12s: -0.56\% average (stability-focused refinement)
\end{itemize}

This trajectory suggests that YOLO11s represented a significant architectural breakthrough for agricultural applications, while YOLO12s focuses on enhancing reliability rather than pushing peak performance boundaries.

\subsection{Hardware Efficiency and Computational Considerations}

While detailed hardware metrics are documented in visualization artifacts, key computational characteristics of YOLO12s include:

\begin{table}[htbp]
\centering
\caption{YOLO12s Computational Efficiency Summary}
\label{tab:yolo12s_efficiency}
\begin{tabular}{lcccc}
\toprule
\textbf{Metric} & \textbf{Coarse1} & \textbf{CropOrWeed2} & \textbf{Fine24} & \textbf{CropsOrWeed9} \\
\midrule
Avg Training Time/Fold & 145 min & 152 min & 163 min & 158 min \\
Peak GPU Memory & 8.2 GB & 9.1 GB & 10.8 GB & 9.7 GB \\
Avg GPU Utilization & 82\% & 84\% & 86\% & 85\% \\
Inference Speed (FPS) & 47 & 45 & 41 & 43 \\
Model Size & 22.3 MB & 22.3 MB & 22.3 MB & 22.3 MB \\
\bottomrule
\end{tabular}
\end{table}

\textbf{Efficiency Highlights:}
\begin{itemize}
    \item \textbf{Consistent resource footprint:} Model size remains constant across variants, simplifying deployment planning
    \item \textbf{Real-time capable:} Inference speeds of 41-47 FPS enable real-time agricultural robotics applications
    \item \textbf{Moderate GPU requirements:} Peak memory usage <11 GB enables training on mid-range GPUs (RTX 3080, V100)
    \item \textbf{High utilization:} 82-86\% average GPU utilization indicates efficient hardware exploitation
\end{itemize}

\subsection{Future Research Directions Informed by YOLO12s Results}

The YOLO12s experimental campaign suggests several promising research avenues:

\begin{enumerate}
    \item \textbf{Stability-Performance Optimization:}
    \begin{itemize}
        \item Investigate techniques to achieve YOLO11s peak performance with YOLO12s stability
        \item Explore ensemble methods combining both architectures for optimal performance-reliability trade-offs
        \item Develop dynamic architecture selection based on per-sample confidence metrics
    \end{itemize}
    
    \item \textbf{Architecture-Task Alignment:}
    \begin{itemize}
        \item Analyze which architectural components contribute to YOLO12s's enhanced stability
        \item Identify task-specific architectural optimizations for agricultural applications
        \item Develop hybrid architectures leveraging strengths of both YOLO11s and YOLO12s
    \end{itemize}
    
    \item \textbf{Fine-Grained Classification Enhancement:}
    \begin{itemize}
        \item Given YOLO12s's superior stability on Fine24, investigate specialized training protocols
        \item Explore hierarchical classification approaches leveraging YOLO12s's consistency
        \item Develop uncertainty quantification methods to flag low-confidence fine-grained predictions
    \end{itemize}
    
    \item \textbf{Deployment Optimization:}
    \begin{itemize}
        \item Quantize YOLO12s for edge deployment while preserving stability advantages
        \item Develop adaptive inference strategies that leverage YOLO12s's predictable behavior
        \item Create deployment frameworks with automatic YOLO11s/YOLO12s selection based on field conditions
    \end{itemize}
\end{enumerate}

\subsection{Conclusions from YOLO12s Analysis}

The comprehensive YOLO12s experimental campaign yields several critical conclusions:

\begin{enumerate}
    \item \textbf{Stability-Performance Trade-off Identified:} YOLO12s sacrifices minimal mean performance (0.56\%) for substantial stability gains (13.1\% CV reduction), representing a valuable architectural design choice for production deployments.
    
    \item \textbf{Task-Specific Architectural Advantages:} YOLO12s excels particularly on complex tasks (Fine24, Coarse1) where consistency is critical, while YOLO11s maintains slight advantages on intermediate complexity tasks (CropOrWeed2).
    
    \item \textbf{Maturation of Agricultural YOLO Architectures:} Both YOLO11s and YOLO12s achieve >70\% mAP50-95 on CropsOrWeed9, suggesting agricultural computer vision has reached practical deployment thresholds.
    
    \item \textbf{Deployment Readiness:} The narrow performance gap between YOLO11s and YOLO12s, combined with YOLO12s's superior stability, indicates both architectures are production-ready with clear use-case differentiation.
\end{enumerate}

These findings establish YOLO12s as a \textit{reliability-optimized} alternative to YOLO11s, expanding deployment options for agricultural robotics applications with varying performance-consistency requirements.
